% GENERATED FILE. Edit canonical lessons, then run npm run generate:books.
% canonical-source-hash: fnv1a64:6aef41bf647c83e7
% canonical-lessons: KA-C258-majjige, KA-C258-benne, KA-C258-payasa, KA-C258-biskattu, KA-C258-bred

\chapter{Buttermilk, Butter, Sweet milk pudding, Biscuit, Bread}
\label{ch:ka-buttermilk-butter-sweet-milk-pudding-biscuit-bread}
\hlchaptermodality{258}

\begin{chapteropening}
\textbf{By the end of this chapter:} \emph{I can say buttermilk, butter, sweet milk pudding, a biscuit and bread.}

Complete the last lesson of chapter 258: I can say buttermilk, butter, sweet milk pudding, a biscuit and bread.
\end{chapteropening}

\section[\texorpdfstring{\kn{ಮಜ್ಜಿಗೆ} (majjige)}{ಮಜ್ಜಿಗೆ (majjige)}]{\kn{ಮಜ್ಜಿಗೆ} (majjige) — buttermilk}

\label{lesson:KA-C258-majjige}



\begin{quote}
Before the new one: say the Kannada for lentils, then the Kannada for tamarind.
\end{quote}

\subsection*{You'll want to know: \kn{ಮಜ್ಜಿಗೆ}}
\textbf{\kn{ಮಜ್ಜಿಗೆ}} — \emph{majjige} — \textquotedblleft{}buttermilk\textquotedblright{}.

\begin{grammarlens}[title={What you've built}]
One more word for this chapter.
\end{grammarlens}

\subsection*{Your turn}
\begin{itemize}
\raggedright
  \item \emph{Say it:} \emph{majjige}
  \item \emph{Say it:} \emph{majjige}, once more
  \item \emph{Say it:} say \emph{majjige}
  \item \emph{Recall:} say the Kannada for lentils, then the Kannada for tamarind, then say \emph{majjige} again
\end{itemize}

\begin{tcolorbox}[breakable,colback=teal!4,colframe=teal!35!black,title={Before you move on}]
What does \kn{ಮಜ್ಜಿಗೆ} mean? (\textquotedblleft{}Buttermilk\textquotedblright{}.) Say it once more.
\end{tcolorbox}
\section[\texorpdfstring{\kn{ಬೆಣ್ಣೆ} (beṇṇe)}{ಬೆಣ್ಣೆ (beṇṇe)}]{\kn{ಬೆಣ್ಣೆ} (beṇṇe) — butter}

\label{lesson:KA-C258-benne}



\begin{quote}
Before the new one: say the Kannada for tamarind, then the Kannada for buttermilk.
\end{quote}

\subsection*{You'll want to know: \kn{ಬೆಣ್ಣೆ}}
\textbf{\kn{ಬೆಣ್ಣೆ}} — \emph{beṇṇe} — \textquotedblleft{}butter\textquotedblright{}.

\begin{grammarlens}[title={What you've built}]
One more word for this chapter.
\end{grammarlens}

\subsection*{Your turn}
\begin{itemize}
\raggedright
  \item \emph{Say it:} \emph{beṇṇe}
  \item \emph{Say it:} \emph{beṇṇe}, once more
  \item \emph{Say it:} say \emph{beṇṇe}
  \item \emph{Recall:} say the Kannada for tamarind, then the Kannada for buttermilk, then say \emph{beṇṇe} again
\end{itemize}

\begin{tcolorbox}[breakable,colback=teal!4,colframe=teal!35!black,title={Before you move on}]
What does \kn{ಬೆಣ್ಣೆ} mean? (\textquotedblleft{}Butter\textquotedblright{}.) Say it once more.
\end{tcolorbox}
\section[\texorpdfstring{\kn{ಪಾಯಸ} (pāyasa)}{ಪಾಯಸ (pāyasa)}]{\kn{ಪಾಯಸ} (pāyasa) — sweet milk pudding, payasa}

\label{lesson:KA-C258-payasa}



\begin{quote}
Before the new one: say the Kannada for buttermilk, then the Kannada for butter.
\end{quote}

\subsection*{You'll want to know: \kn{ಪಾಯಸ}}
\textbf{\kn{ಪಾಯಸ}} — \emph{pāyasa} — \textquotedblleft{}sweet milk pudding, payasa\textquotedblright{}.

\begin{grammarlens}[title={What you've built}]
One more word for this chapter.
\end{grammarlens}

\subsection*{Your turn}
\begin{itemize}
\raggedright
  \item \emph{Say it:} \emph{pāyasa}
  \item \emph{Say it:} \emph{pāyasa}, once more
  \item \emph{Say it:} say \emph{pāyasa}
  \item \emph{Recall:} say the Kannada for buttermilk, then the Kannada for butter, then say \emph{pāyasa} again
\end{itemize}

\begin{tcolorbox}[breakable,colback=teal!4,colframe=teal!35!black,title={Before you move on}]
What does \kn{ಪಾಯಸ} mean? (\textquotedblleft{}Sweet milk pudding, payasa\textquotedblright{}.) Say it once more.
\end{tcolorbox}
\section[\texorpdfstring{\kn{ಬಿಸ್ಕತ್ತು} (biskattu)}{ಬಿಸ್ಕತ್ತು (biskattu)}]{\kn{ಬಿಸ್ಕತ್ತು} (biskattu) — a biscuit}

\label{lesson:KA-C258-biskattu}



\begin{quote}
Before the new one: say the Kannada for butter, then the Kannada for sweet milk pudding.
\end{quote}

\subsection*{You'll want to know: \kn{ಬಿಸ್ಕತ್ತು}}
\textbf{\kn{ಬಿಸ್ಕತ್ತು}} — \emph{biskattu} — \textquotedblleft{}a biscuit\textquotedblright{}.

\begin{grammarlens}[title={What you've built}]
One more word for this chapter.
\end{grammarlens}

\subsection*{Your turn}
\begin{itemize}
\raggedright
  \item \emph{Say it:} \emph{biskattu}
  \item \emph{Say it:} \emph{biskattu}, once more
  \item \emph{Say it:} say \emph{biskattu}
  \item \emph{Recall:} say the Kannada for butter, then the Kannada for sweet milk pudding, then say \emph{biskattu} again
\end{itemize}

\begin{tcolorbox}[breakable,colback=teal!4,colframe=teal!35!black,title={Before you move on}]
What does \kn{ಬಿಸ್ಕತ್ತು} mean? (\textquotedblleft{}A biscuit\textquotedblright{}.) Say it once more.
\end{tcolorbox}
\section[\texorpdfstring{\kn{ಬ್ರೆಡ್} (breḍ)}{ಬ್ರೆಡ್ (breḍ)}]{\kn{ಬ್ರೆಡ್} (breḍ) — bread}

\label{lesson:KA-C258-bred}



\begin{quote}
Before the new one: say the Kannada for sweet milk pudding, then the Kannada for a biscuit.
\end{quote}

\subsection*{You'll want to know: \kn{ಬ್ರೆಡ್}}
\textbf{\kn{ಬ್ರೆಡ್}} — \emph{breḍ} — \textquotedblleft{}bread\textquotedblright{}.

\begin{grammarlens}[title={What you've built}]
One more word for this chapter.
\end{grammarlens}

\subsection*{Your turn}
\begin{itemize}
\raggedright
  \item \emph{Say it:} \emph{breḍ}
  \item \emph{Say it:} \emph{breḍ}, once more
  \item \emph{Say it:} say \emph{breḍ}
  \item \emph{Recall:} say the Kannada for sweet milk pudding, then the Kannada for a biscuit, then say \emph{breḍ} again
\end{itemize}

\begin{tcolorbox}[breakable,colback=teal!4,colframe=teal!35!black,title={Before you move on}]
What does \kn{ಬ್ರೆಡ್} mean? (\textquotedblleft{}Bread\textquotedblright{}.) Say it once more.
\end{tcolorbox}
